% page 169 of the facsimile (image p-168 of the scan)
% block 165.1 x 237.7 mm ; left margin 36.9 mm
\pgc{15.240mm}{0.000mm}{
\l{\cel{56}159}
\l{}
\l{}
\l{\sou{fardo}. I. Kompozajo quan onu aplikas sur la pelo por plubeligar}
\l{\cel{3}la karnaciono. - II. (\sou{metaf}.)Aspekto versemblanta qua kolo-}
\l{\cel{3}rizas o travestias la fakti. - F.}
\l{}
\l{\sou{farino}.\cel{2}Substanco nutriva di ula grani de cereali, pulverigita.}
\l{\cel{3}- EFIS.}
\l{}
\l{\sou{faringo}. (\sou{anat.)}\cel{3}La dopa parto di la boko.}
\l{}
\l{\sou{faringito}. (\sou{patol.)}\cel{2}Inflamuro di la faringo.}
\l{}
\l{\sou{farizeo}. La homo qua agas la gesti religiala, ma di qua la pieso}
\l{\cel{3}esas plu samblanta kam reala. - DEFIRS.}
\l{}
\l{\sou{farm}ar. (\sou{trans.}) Signatar la konvenciono per qua proprietero}
\l{\cel{3}lugas ter-domeno ad ulu qua explotos olu po preco yarala.-DEF.}
\l{}
\l{\sou{farmacio}. La arto preparar la medikamenti. - DEFIRS.}
\l{}
\l{\sou{farmakopeo}. Libro qua traktas la arto preparar la medikamenti.-}
\l{\cel{3}DEFIRS.}
\l{}
\l{\sou{farso}.\cel{2}I. Hachajo de karni spicizita ye qua onu garnisas la}
\l{\cel{3}parto interna di pultro, di pasteto, e c. - II. Mikra teatrajo}
\l{\cel{3}bufonala. - III. Kozo bufonala quan onu dicas od agas.-DEFIRS.}
\l{}
\l{\sou{fasado}. Parto avana di konstrukturo, ube esas la enireyo pre-}
\l{\cel{3}cipua. - DEFIS.}
\l{}
\l{\sou{\textquotedbl{}fascia\textquotedbl{}}. (\sou{histologio}). Termino per qua indikesas formacuri}
\l{\cel{3}aponevrosala, veli sive fibra sive celulara, qui kovras}
\l{\cel{3}od envelopas muskuli o regioni.}
\l{}
\l{\sou{fashino}. Branchari asemblita a faski, quin onu uzas por plenigar}
\l{\cel{3}la fosati di fortifikajo, shirm-apogar kanono-baterii, impe-}
\l{\cel{3}dar la krulo di tero, e c. - DEFIRS.}
\l{}
\l{fashismo. Sistemo politikala, qua havas karaktero nacionalista,}
\l{\cel{3}antisocialista, antiparlamentista, instaurita (en 1922)}
\l{\cel{3}en Italia, da Benito Mussolini. - DEFIRS.}
\l{}
\l{fasiacar. (\sou{netrans.})\cel{2}(\sou{bot.)}\cel{3}Deformesar teratologie per aplas-}
\l{\cel{3}to (stipi, rami, pedunkli). - EF.}
\l{}
\l{fasko. I. Asemblajo de ula quanto de kozi simila, longa, inter-}
\l{\cel{3}ligita. - II. (\sou{optiko).}\cel{1}Ensemblo ek radii o de lumo-pinseli}
\l{\cel{3}qui emanas de un fonto. - FIS.}
\l{}
\l{fasonar. (\sou{trans}.) Laborar ula materio por fabrikar objekto de-}
\l{\cel{3}finita (fasono di vesto, di juvelo). - DEFR.}
\l{}
\l{fastar. (\sou{netrans., dum).}\cel{1}Abstenar plu o min komplete nutrivi. DE.}
}
